\section{Zadanie 1}
Wyprowadź z aksjomatów kolejne własności przestrzeni probabilistycznej

\subsection*{Aksjomaty miary probabilistycznej}
\begin{itemize}
	\item \(P \colon \mathcal F \to [0, \infty)\)
	\item \(P(\Omega) = 1\)
	\item jeśli \(\{A_n\}_{n\in\mathbb N}\) jest (przeliczalną!) rodziną
	      zbiorów parami rozłącznych,\\ to
	      \(P(\bigcup_{n\in\mathbb N} A_n) = \sum_{n\in\mathbb N} P(A_n)\)
\end{itemize}

\subsection*{Rozwiązanie}
\begin{enumerate}
	\item[(i)]
	      \(P(\emptyset) = 0\)\\
	      \(\{\Omega, \emptyset, \emptyset, \emptyset, \emptyset, \ldots\}\)
	      jest przeliczalną rodziną zbiorów parami rozłącznych\\
	      Z aksjomatów:\\
	      \(
	      P(\Omega \cup \emptyset \cup \emptyset \cup \ldots) =
	      P(\Omega) + P(\emptyset) + P(\emptyset) + \ldots
	      \)\\
	      \(
	      P(\Omega) =
	      P(\Omega) + P(\emptyset) + P(\emptyset) + \ldots
	      \)\\
	      \(
	      1 = 1 + P(\emptyset) + P(\emptyset) + \ldots
	      \)\\
	      \(
	      0 = \sum_{n\in\mathbb N} P(\emptyset)
	      \)\\
	      \(
	      P(\emptyset) = 0
	      \)\\
	      Korzystając z tej tożsamości będziemy od tej pory mogli używać
	      aksjomatu 3. w wersji skończonej - nieskończenie wiele zbiorów
	      pustych nie zmieni wyniku.

	\item[(ii)]
	      \(P(\Omega\setminus A) = 1 - P(A)\)\\
	      \(\{(\Omega\setminus A), A\}\) to skończona rodzina zbiorów parami rozłącznych.
	      W zgodzie z aksjomatem powinniśmy w niej uwzględnić jeszcze nieskończenie
	      wiele zbiorów pustych\\
	      \(P((\Omega\setminus A) \cup A) = P(\Omega\setminus A) + P(A)\)\\
	      \(P(\Omega) = P(\Omega\setminus A) + P(A)\)\\
	      \(P(\Omega\setminus A) = P(\Omega) - P(A)\)\\
	      \(P(\Omega\setminus A) = 1 - P(A)\)

	\item[(iii)]
	      jeśli \(A \subseteq B\), to \(P(B\setminus A) = P(B) - P(A)\)\\
	      \((B\setminus A), A\) są rozłączne:\\
	      \(P((B\setminus A)\cup A) = P(B\setminus A) + P(A)\)\\
	      \(P(B) = P(B\setminus A) + P(A)\)\\
	      \(P(B\setminus A) = P(B) - P(A)\)

	\item[(iv)]
	      jeśli \(A \subseteq B\), to \(P(A) \leqslant P(B)\)\\
	      korzystając z (iii): \(P(B) = P(A) + P(B\setminus A)\)\\
	      z aksjomatu \(P\) przyjmuje wartości nieujemne:
	      \(P(B\setminus A) \geqslant 0\)\\
	      \(P(B) = P(A) + P(B\setminus A) \geqslant P(A)\)

	\item[(v)]
	      \(P(A) \leqslant 1\)\\
	      korzystamy z (ii): \(P(A) = 1 - P(\Omega\setminus A)\)\\
	      \(P\) przyjmuje wartości nieujemne: \(P(\Omega\setminus A) \geqslant 0\)\\
	      \(P(A) = 1 - P(\Omega\setminus A) \leqslant 1\)

	\item[(vi)]
	      \(P(A\cup B) = P(A) + P(B) - P(A\cap B)\)\\
	      Zauważmy \(P(B) = P(B\setminus A) + P(A\cap B)\)\\
	      \(P(A\cup B) = P(A) + P(B\setminus A) = P(A) + P(B) - P(A\cap B)\)
\end{enumerate}
